% META: {"id": "TCOMPL-BAYES-CR-010", "chapitre": "TCOMPL-INFERENCE-BAYESIENNE", "type_objet": "cours", "capacites_codes": ["C1", "C2", "C3"], "sources_inspiration": [], "mode_creation": "ex_nihilo", "status": "generated"}

\section{La formule de Bayes}

% ============================================================
% STRATE 1 — Rappel proba conditionnelle, formule de Bayes
% ============================================================

\definition[1]{
Pour deux événements $A$ et $B$ avec $P(A)>0$, la \textbf{probabilité conditionnelle} de $B$ sachant $A$ est $P_A(B) = \dfrac{P(A\cap B)}{P(A)}$.
}

\margeAppui{
On appelle $P(A)$ la \textbf{probabilité a priori} de $A$ (avant toute observation), et $P_B(A)$ la \textbf{probabilité a posteriori} de $A$ (après avoir observé $B$).
}

\theoreme[Formule de Bayes]{
Pour $A$ et $B$ deux événements de probabilité non nulle :
\[
  P_B(A) = \frac{P_A(B) \times P(A)}{P(B)}.
\]
}

\demonstration{
Par définition de la probabilité conditionnelle, $P(A\cap B) = P_A(B)\times P(A)$ et aussi $P(A\cap B) = P_B(A)\times P(B)$ (les deux expressions désignent la même probabilité $P(A\cap B)$, calculée dans les deux sens). En égalant : $P_B(A)\times P(B) = P_A(B)\times P(A)$, d'où le résultat en divisant par $P(B)$ (non nul).
}

% ============================================================
% STRATE 2 — Distinction PB(A) vs PA(B)
% ============================================================

\erreurFrequente{
\textbf{Confondre $P_B(A)$ et $P_A(B)$.} Ces deux quantités sont en général \textbf{très différentes} : $P_A(B)$ est la probabilité de $B$ sachant $A$ (par exemple, la probabilité qu'un test soit positif sachant que le patient est malade), alors que $P_B(A)$ est la probabilité de $A$ sachant $B$ (la probabilité que le patient soit malade sachant que le test est positif). C'est précisément la formule de Bayes qui permet de passer de l'une à l'autre.
}

\exemple{
Dans une population, $1\%$ des personnes sont porteuses d'une maladie ($P(A)=0{,}01$). Un test détecte la maladie avec une probabilité de $99\%$ quand elle est présente ($P_A(B)=0{,}99$), et donne un résultat positif à tort dans $2\%$ des cas chez les personnes saines. On calcule $P(B) = P(A)P_A(B) + P(\bar{A})P_{\bar{A}}(B) = 0{,}01\times0{,}99 + 0{,}99\times0{,}02 = 0{,}0297$.

La probabilité qu'une personne testée positive soit réellement malade est :
\[
  P_B(A) = \frac{P_A(B)\times P(A)}{P(B)} = \frac{0{,}99\times0{,}01}{0{,}0297} \approx 0{,}333.
\]
}

% BEGIN-VERIFY
% from sympy import Rational
% PA = Rational(1, 100)
% PA_B = Rational(99, 100)
% PnotA_B = Rational(2, 100)
% PB = PA*PA_B + (1-PA)*PnotA_B
% assert PB == Rational(297, 10000)
% PB_A = PA_B * PA / PB
% assert abs(float(PB_A) - 0.333) < 0.001
% END-VERIFY

\propriete[Le paradoxe du test fiable]{
Même avec un test fiable à $99\%$, la probabilité $P_B(A)$ d'être réellement malade sachant le test positif ne vaut que $33\%$ environ dans l'exemple ci-dessus : cela s'explique par la \textbf{rareté} de la maladie ($1\%$ de la population), qui fait que le nombre de \textbf{faux positifs} (personnes saines testées à tort positives) devient important en valeur absolue, même avec un faible taux d'erreur individuel.
}
